\section{Deployment \& Infrastructure}

\subsection{Host Configuration}
All components (Mosquitto, Telegraf, \texttt{daq2-ingest}, InfluxDB, and Grafana) are co-deployed as Docker containers on a single historian server running Ubuntu Server 24.04 LTS on the lab LAN. Co-deployment is justified at the two-node testbed scale. At the full rate, CPU is the limiting resource: approximately 1.5 cores of InfluxDB and 0.6 cores of ingestion per node (Section~\ref{sec:fullrate}), so the host should be sized to the number of nodes.

\begin{table}[H]
\centering
\caption{Component Deployment Summary}
\begin{tabular}{llll}
\hline
\textbf{Component} & \textbf{Image} & \textbf{Host port} & \textbf{Purpose} \\
\hline
Mosquitto & \texttt{eclipse-mosquitto:2}     & \texttt{127.0.0.1:1883}$^\dagger$ & MQTT broker \\
InfluxDB  & \texttt{influxdb:2}              & \texttt{127.0.0.1:8086} & Time-series database \\
Telegraf  & \texttt{telegraf:latest}         & none                    & MQTT to InfluxDB bridge (JSON) \\
\texttt{daq2-ingest} & local build      & none                    & MQTT to InfluxDB bridge (full rate) \\
Grafana   & \texttt{grafana/grafana:latest}  & \texttt{0.0.0.0:3000}   & Dashboards \\
\hline
\multicolumn{4}{l}{\small $^\dagger$ Changed to \texttt{1883:1883} when real DAQ nodes are connected.}
\end{tabular}
\end{table}

\subsection{Docker Compose}
\label{sec:compose}
All components are managed via a single Docker Compose file. This ensures reproducible deployment, simplifies inter-service networking, and allows the full stack to be brought up or torn down with a single command, appropriate for a research testbed where the environment may need to be rebuilt. The as-built \texttt{docker-compose.yml} is:

\begin{lstlisting}[language=toml]
services:
  mosquitto:
    image: eclipse-mosquitto:2
    container_name: mosquitto
    ports:
      - "127.0.0.1:1883:1883"
    volumes:
      - ./mosquitto/config:/mosquitto/config
      - ./mosquitto/data:/mosquitto/data
      - ./mosquitto/log:/mosquitto/log
    restart: unless-stopped

  influxdb:
    image: influxdb:2
    container_name: influxdb
    ports:
      - "127.0.0.1:8086:8086"
    volumes:
      - influxdb-data:/var/lib/influxdb2
      - ./influxdb/init:/docker-entrypoint-initdb.d
    environment:
      - DOCKER_INFLUXDB_INIT_MODE=setup
      - DOCKER_INFLUXDB_INIT_USERNAME=${INFLUXDB_ADMIN_USER}
      - DOCKER_INFLUXDB_INIT_PASSWORD=${INFLUXDB_ADMIN_PASSWORD}
      - DOCKER_INFLUXDB_INIT_ORG=qut-microgrid
      - DOCKER_INFLUXDB_INIT_BUCKET=telemetry
      - DOCKER_INFLUXDB_INIT_RETENTION=90d
      - DOCKER_INFLUXDB_INIT_ADMIN_TOKEN=${INFLUXDB_TOKEN}
      # Compactions: by default up to half the CPU threads run at once. Several full compactions
      # of hourly waveform shards together pushed InfluxDB to ~11 GB, swapped the host and stalled
      # writes for minutes (lab soak, 14 Oct 2026). Two at a time keeps up with two nodes.
      - INFLUXD_STORAGE_MAX_CONCURRENT_COMPACTIONS=2
      # Soft heap limit: the Go runtime collects harder near it instead of growing into swap.
      # Not a hard cap (a Docker mem_limit would kill InfluxDB instead). Size it below host RAM.
      - GOMEMLIMIT=${INFLUXDB_MEMORY_LIMIT:-8GiB}
    restart: unless-stopped

  telegraf:
    image: telegraf:latest
    container_name: telegraf
    depends_on:
      - mosquitto
      - influxdb
    volumes:
      - ./telegraf/telegraf.conf:/etc/telegraf/telegraf.conf:ro
    environment:
      - INFLUXDB_TOKEN=${INFLUXDB_TOKEN}
    restart: unless-stopped

  daq2-ingest:
    build: ./daq2-ingest
    container_name: daq2-ingest
    depends_on:
      - mosquitto
      - influxdb
    environment:
      - INFLUXDB_TOKEN=${INFLUXDB_TOKEN}
      # Raw full-rate retention, only applied when the bucket is first created
      - RAW_RETENTION=${RAW_RETENTION:-7d}
    restart: unless-stopped

  grafana:
    image: grafana/grafana:latest
    container_name: grafana
    ports:
      - "3000:3000"
    volumes:
      - grafana-data:/var/lib/grafana
      - ./grafana/provisioning:/etc/grafana/provisioning
    environment:
      - GF_SECURITY_ADMIN_PASSWORD=${GRAFANA_ADMIN_PASSWORD}
      # Read by the provisioned data source (grafana/provisioning/datasources/influxdb.yml)
      - INFLUXDB_TOKEN=${INFLUXDB_TOKEN}
      - GF_AUTH_ANONYMOUS_ENABLED=false
      - GF_DASHBOARDS_MIN_REFRESH_INTERVAL=1s
    depends_on:
      - influxdb
    restart: unless-stopped

volumes:
  influxdb-data:
  grafana-data:
\end{lstlisting}

The InfluxDB and Mosquitto ports are bound to loopback by default; Grafana is published on all interfaces for dashboard access from the lab LAN (Section~\ref{sec:security}). Sensitive values are passed via environment variables sourced from a \texttt{.env} file that is excluded from version control.

\subsection{Compaction and Memory}
\label{sec:compaction}
The \texttt{waveform} bucket uses one-hour shard groups. About four hours after a shard stops receiving writes, InfluxDB rewrites it in a full compaction, which on the lab server's HDD takes 20--25 minutes for one node's hour (about 2.5~GB). By default InfluxDB runs up to half as many compactions at once as there are CPU threads. In the lab soak, three or four compactions overlapped twice: InfluxDB grew to about 11~GB, the host swapped, writes stalled for several minutes, and the node lost packets because it could not publish.

The Compose file therefore sets \texttt{INFLUXD\_STORAGE\_MAX\_CONCURRENT\_COMPACTIONS=2}, so further compactions queue instead of running together, and \texttt{GOMEMLIMIT} (default 8~GiB, set with \texttt{INFLUXDB\_MEMORY\_LIMIT} in \texttt{.env}), a soft limit at which the Go runtime collects garbage harder instead of growing into swap. \texttt{GOMEMLIMIT} should sit a few gigabytes below the host's RAM. Two concurrent compactions suit up to two nodes; if compactions fall behind with more nodes (\texttt{TSM compaction} lines in \texttt{docker logs influxdb} queueing for longer than an hour), raise the limit or move InfluxDB to faster storage.

\subsection{Persistent Volumes}
InfluxDB and Grafana state are stored in named Docker volumes. Mosquitto data and logs are bind-mounted to host directories for visibility. The Telegraf configuration is bind-mounted read-only, and Grafana's data source and dashboard are provisioned from \path{grafana/provisioning/}. \texttt{daq2-ingest} is stateless: its image is built from the repository, and it creates its buckets on first start.

\subsection{Environment File}
A \texttt{.env} file in the Compose project root holds all secrets:

\begin{lstlisting}[language=toml]
INFLUXDB_ADMIN_USER=admin
INFLUXDB_ADMIN_PASSWORD=<redacted>   # minimum 8 characters
INFLUXDB_TOKEN=<redacted>
GRAFANA_ADMIN_PASSWORD=<redacted>
RAW_RETENTION=7d                     # optional: full-rate retention
INFLUXDB_MEMORY_LIMIT=8GiB           # optional: InfluxDB soft memory limit
\end{lstlisting}

This file must not be committed to version control. Its values are applied only on the first start of the containers (\texttt{RAW\_RETENTION} only when the \texttt{waveform} bucket is first created); changing them later requires resetting the volumes. \texttt{INFLUXDB\_MEMORY\_LIMIT} is the exception: it is applied whenever the InfluxDB container is created (Section~\ref{sec:compaction}). Step-by-step setup instructions are given in the repository \texttt{README.md}.
